\section{Wstęp}

W dzisiejszych czasach dostęp do sztucznej inteligencji stał się powszechny, a jej rozwój odbywa się w błyskawicznym tempie. Zawładnęła ona takimi dziedzinami, jak media społecznościowe, zakupy internetowe, medycyna, transport i nawigacja, bankowość, gry, analityka czy biznes \cite{tableau_software}. Coraz silniej przenika do naszego codziennego życia i staje się jego integralnym elementem. W 2025 roku technologii AI, w takich sektorach jak komunikacja czy informacja, używało 19,95\% wszystkich europejskich przedsiębiorstw i 55,03\% dużych europejskich przedsiębiorstw \cite{eu_dane_z_firm}. Przykładem mogą tutaj być samochody autonomiczne, które są pojazdami poruszającymi się bez udziału człowieka, wygenerowane zdjęcia czy artykuły. Ta niezwykle dynamiczna dziedzina budzi w wielu z nas, bo aż w trzech na pięć osób, niepokój przeplatany z ekscytacją. Pragniemy postrzegać ją jako bezstronną i działającą w sposób wpierający użytkownika, ale jesteśmy wobec niej nieufni, albowiem jak przekonać się do czegoś, co nawet nie wiemy jak działa? Czy możemy powierzyć jej decyzje wpływające na nasze życia, bezpieczeństwo, finanse i zdrowie? Skąd możemy być pewni, że kieruje się naszym dobrem oraz wymienionymi wcześniej wartościami? Czy powinniśmy na ślepo dążyć za wizją postępu bez sprawdzenia, czy nasze postępowanie jest moralne? \cite{uni_queen} 

Bez odpowiedzi na owe dylematy nie powinniśmy powierzać sztucznej inteligencji naszego życia czy zdrowia w tak krytycznych miejscach jak wspomniana wcześniej medycyna, transport czy biznes. Potwierdzają to obawy przedstawione przez przedsiębiorstwa, które  wskazywały m.in. na brak odpowiedniej wiedzy eksperckiej oraz problem dotyczący naruszenia ochrony danych i prywatności. W jaki więc sposób można zbadać systemy wykorzystujące AI? Jak zwiększyć ich niezawodność, zyskać do nich zaufanie i potwierdzić, że~nie stanowią one dla nas zagrożenia? 

\begin{figure}[H]
\centering
\includegraphics[width=0.8\linewidth]{Enterprises_that_have_ever_considered_using_any_of_the_AI_technologies_by_reason_for_not_using_EU_2025.png}
\caption{\label{fig:eu_enterpraises_plot}Przyczyny, z powodu których przedsiębiorstwa rozważające wykorzystanie sztucznej inteligencji w swoich systemach jej nie używały \cite{eu_dane_z_firm}}
\end{figure}

XAI, czyli wyjaśnialna sztuczna inteligencja (ang. \textit{Explainable Artificial Intelligence}) to dziedzina składająca się z różnych procedur i technik mających na celu wyjaśnienie modeli AI. Robi to za pomocą algorytmów, które starają się stworzyć odpowiedź możliwą do zrozumienia dla ludzkich odbiorców \cite{gfg_xai_definition}. Nie w każdym z przypadków użytkownik musi znać szczegółowe przyczyny podjęcia decyzji przez sztuczną inteligencję. Przykładowo, podczas poruszania się samochodem autonomicznym dla pasażera nie jest istotne to, dlaczego konkretny obiekt został zaklasyfikowany w określony sposób, a to, czy auto zauważyło przeszkodę i się zatrzymało. XAI nie musi być używane wyłącznie do interpretacji wyników, jak~powinno to mieć miejsce w medycynie. Może zostać użyte także w trakcie produkcji modelu, aby wykryć błędne założenia budowanej struktury, jej zachowania dyskryminacyjne, czy też błędy ludzkie \cite{molnar2025}. 

Celem pracy jest opracowanie i porównanie metod interpretacji modeli uczenia maszynowego stosowanych do klasyfikacji danych różnych dziedzin w kontekście zwiększenia przejrzystości i zaufania do decyzji systemów AI wspomagających te dziedziny. W~tym~celu wyselekcjonowano cztery zbiory danych, z czego 2 składają się z obrazów, do których zostaną zbudowane lub wykorzystane już istniejące modele, aby następnie poddać je analizie, żeby sprawdzić, w jaki sposób przewidują wynik końcowy. Zamiarem jest zastosowanie czterech różnych metody XAI, tj. LIME, SHAP, Grad–CAM oraz PFI, które~będą porównywane pod kątem przejrzystości wyników, zrozumiałości dla~użytkownika końcowego (np. dla danych medycznych – lekarza) oraz wpływu na~dokładność modelu. Przyjęto tezę, że obecnie sztuczna inteligencja wymaga interwencji człowieka, który powinien sprawdzać i kontrolować jej działanie. Opinia ta znajduje uzasadnienie w~istnieniu tzw. halucynacji – zdawałoby się brzmiących bezbłędnie, nie budzących wątpliwości informacji, np. w postaci danych statystycznych, definicji, źródeł, które tak naprawdę nie mają miejsca w świecie rzeczywistym (nie istnieją), prawdopodobnie będąc stworzone na~rzecz omawianego problemu celem wygenerowania zadowalającej odpowiedzi za wszelką cenę \cite{halucination_definition}; lub tym, że systemy AI mają tendencję do kierowania się błędnymi korelacjami, czyli~współzależnościami w danych \cite{corellation_definition}.

Dokument składa się z 6 rozdziałów. Drugi, zatytułowany ,,Problem klasyfikacji'', wprowadza czytelnika w dziedzinę problemu oraz przedstawia użyte metryki i modele. W trzecim, ,, Metody XAI'', omówiono zastosowane do wyjaśniania metody, w tym ich algorytmy, podstawy teoretyczne, wady i zalety. Rozdział czwarty ,,Opis implementacji'' składa się z najważniejszych elementów implementacji, sposobu przeprowadzania eksperymentów, charakterystyk zbiorów i sposobów ich przygotowania. ,,Eksperymenty'' prezentują otrzymane w ramach badań wyniki wraz z ich opisami i podsumowaniami. Ostatni rozdział zawiera streszczenie tego, co udało się zrealizować, wnioski ogólne, a także kierunki dalszego rozwoju pracy.
